% =====================================================================
\chapter{Телеметрия}
\label{ch:telemetry}
% =====================================================================
\chapterintro
  {как модель рассказывает пульту о себе: качество связи, заряд батареи, координаты;
   как найти датчики, настроить тревоги и экраны}
  {модель с приёмником ELRS (или другим с телеметрией); для батареи --- полётный контроллер}

\section{Обратная связь}

\begin{simple}
Телеметрия --- это «письма от модели». Приёмник рассказывает пульту, хорошо ли он тебя
слышит, а полётный контроллер добавляет: сколько вольт в батарее, какой ток, где модель по GPS.
Пульт показывает это на экране и может сказать голосом.
\end{simple}

\begin{figure}[H]
\centering
\begin{tikzpicture}[every node/.style={font=\sffamily\small}]
  \pic[transform shape,scale=1.1] at (0,1.2) {battery};
  \node[lbl] at (0,0.6) {батарея};
  \node[draw=black!50,fill=white,rounded corners=2pt,font=\sffamily\scriptsize] at (0,-0.4) {GPS};
  \node[lbl] at (0,-0.9) {спутники};
  \pic[transform shape,scale=1.1] at (2.4,0.4) {fcboard};
  \node[lbl,align=center] at (2.4,-0.6) {полётный\\контроллер};
  \draw[arr] (0.8,1.1) -- (1.7,0.6);
  \draw[arr] (0.5,-0.4) -- (1.7,0.2);
  \pic[transform shape,scale=1.2] at (5.0,0.4) {receiver};
  \node[lbl] at (5.0,-0.3) {приёмник};
  \draw[arr] (3.2,0.4) -- (4.1,0.4);
  \draw[arr,dashed,etx,line width=1.2pt] (6.4,0.4) -- node[above,lbl]{по радио} (8.3,0.4);
  \pic[transform shape,scale=1.1] at (9.5,0.4) {minitx};
  \node[draw=black!40,fill=lcdbg,font=\ttfamily\scriptsize,align=left,anchor=west] at (10.8,0.4) {RxBt 15.8V\\RQly 100\%\\Sats 12};
\end{tikzpicture}
\caption{Путь телеметрии: датчики → полётный контроллер → приёмник → пульт}
\end{figure}

\section{Найти датчики}

\begin{enumerate}
  \item Включи модель, дождись связи.
  \item \mpath{MDL\sep Telemetry} → \menu{Discover new sensors}. Пульт начнёт собирать все
        датчики, которые присылает модель.
  \item Через 10 секунд нажми \menu{Stop discovery}.
\end{enumerate}

\begin{center}
\lcdscreen{TELEMETRY}{13/13}{\lr{1RSS}{-54dB}\lr{RQly}{100\%}\lr{TPWR}{100mW}\lr{RxBt}{15.8V}\lr{Curr}{12.3A}\lsel{Discover new sensors}{}}
\end{center}

Мигающее значение или «?» означает, что данные сейчас не приходят.
\menu{Delete all sensors} --- очистить список и начать заново.

\section{Главные датчики ExpressLRS}

\begin{simple}
Представь, что ты разговариваешь с другом через шумную улицу. \textbf{RSSI} --- это
\emph{громкость}: насколько сильно слышен голос. \textbf{LQ} (качество связи) --- это
\emph{сколько слов друг понял}. Можно слышать громко, но через шум понимать мало. Поэтому
главный показатель --- \textbf{LQ}.
\end{simple}

\begin{figure}[H]
\centering
\begin{tikzpicture}[every node/.style={font=\sffamily\small}]
  \fill[warn!60] (0,0) rectangle (5,0.7);
  \fill[yellow!70!orange] (5,0) rectangle (7,0.7);
  \fill[okgreen!60] (7,0) rectangle (10,0.7);
  \draw[black!50] (0,0) rectangle (10,0.7);
  \foreach \x/\l in {0/0,5/50,7/70,10/100}{\draw (\x,0) -- (\x,-0.15) node[below,font=\sffamily\scriptsize]{\l\,\%};}
  \node[font=\sffamily\scriptsize\bfseries,text=white] at (2.5,0.35) {опасно: разворачивайся!};
  \node[font=\sffamily\scriptsize\bfseries] at (6,0.35) {осторожно};
  \node[font=\sffamily\scriptsize\bfseries] at (8.5,0.35) {отлично};
  \node[anchor=east,font=\sffamily\small\bfseries] at (-0.2,0.35) {RQly};
  \draw[etx,line width=1.6pt,-{Latex}] (9.2,1.4) -- (9.2,0.75);
  \node[font=\sffamily\scriptsize,text=etx] at (9.2,1.6) {сейчас 92\,\%};
\end{tikzpicture}
\caption{Качество связи RQly: чем ближе к 100\,\%, тем лучше}
\end{figure}

\begin{center}\small
\renewcommand{\arraystretch}{1.2}
\begin{tabularx}{\linewidth}{@{}L{22mm}Y@{}}
\toprule
\textbf{Датчик} & \textbf{Что показывает} \\
\midrule
\sw{RQly} & качество связи, \% --- \textbf{главный}. Ниже 70\,\% --- пора возвращаться \\
\sw{1RSS}, \sw{2RSS} & уровень сигнала на антеннах приёмника, дБм. Около $-50$ --- рядом, около
  $-100$ и ниже --- на пределе \\
\sw{RSNR} & насколько сигнал сильнее шума \\
\sw{TPWR} & мощность, с которой сейчас передаёт пульт \\
\sw{RxBt} & напряжение батареи модели \\
\sw{Curr}, \sw{Capa} & ток и сколько энергии уже потрачено (мА·ч) \\
\sw{GPS}, \sw{Alt}, \sw{GSpd}, \sw{Sats} & координаты, высота, скорость, число спутников \\
\sw{FM} & режим полётного контроллера (\code{ACRO}, \code{ANGL}, \code{!FS!} --- failsafe) \\
\bottomrule
\end{tabularx}
\end{center}

\section{Тревоги}

На странице \menu{Telemetry} есть пороги для качества связи: \menu{Low alarm} (около 70) и
\menu{Critical alarm} (около 50). Пульт скажет «Signal low» или «Signal critical». Если связь
пропала совсем --- «Telemetry lost», вернулась --- «Telemetry recovered».

Тревогу по батарее сделай сам: логический переключатель \code{a<x} с \sw{RxBt}
(глава~\ref{ch:logic}) + специальная функция \menu{Play track} (глава~\ref{ch:special}).

\begin{tip}
Удобнее следить за напряжением \emph{на одну банку}. В Betaflight для этого есть команда
\code{set report\_cell\_voltage = ON} --- тогда \sw{RxBt} придёт уже в пересчёте на банку,
и порог 3,5\,В подойдёт для любой батареи.
\end{tip}

\section{Экраны телеметрии}

Кнопка \btn{TELE} показывает до четырёх экранов. Что на них --- настраивается в
\menu{Telemetry}, раздел \menu{Screens}:
\begin{itemize}
  \item \menu{Nums} --- таблица чисел;
  \item \menu{Bars} --- полоски с минимумом и максимумом;
  \item \menu{Script} --- экран, нарисованный Lua-скриптом (глава~\ref{ch:lua}).
\end{itemize}

\begin{center}
\lcdscreen{Quad-5in}{TELE 1/2}{\lblank\lr{RxBt\ \ 15.8V}{RQly\ \ 100\%}\lr{Curr\ \ 12.3A}{1RSS\ \ -54}\lr{Capa\ \ 812}{TPWR\ \ 100mW}\lr{FM\ \ \ \ ACRO}{Tmr1\ \ 02:37}}
\end{center}

\section{Сброс}

Минимумы, максимумы и счётчики (например, \sw{Capa}) хранятся до сброса. Удобно сделать
специальную функцию \menu{Reset → Telemetry} на кнопку, или сбросить из меню главного экрана
(долгое \btn{ENTER} → \menu{Reset}).

\begin{tryit}
\begin{enumerate}
  \item Найди датчики своей модели. Сколько их?
  \item Сделай экран \menu{Nums} с \sw{RxBt}, \sw{RQly}, \sw{1RSS}, \sw{TPWR}.
  \item Отойди с моделью (без пропеллеров!) на 30--50 метров, поставив мощность 10\,мВт.
        Как меняются \sw{1RSS} и \sw{RQly}?
\end{enumerate}
\end{tryit}

\begin{summary}
\begin{itemize}
  \item Телеметрия --- данные от модели к пульту.
  \item \menu{Discover new sensors} находит все датчики.
  \item Главный показатель связи --- \sw{RQly}: ниже 70\,\% --- возвращайся.
  \item Тревоги связи --- встроенные, тревогу батареи собираешь сам.
\end{itemize}
\end{summary}
